\chapter{月面のモザイク合成}\label{chap:mosaic}

焦点距離の長い望遠鏡では、月の全体が1つの画面に収まりません。月を\textbf{何か所かに分けて撮り}（それぞれを\textbf{パネル}と呼びます）、
あとで1枚につなぐのが\textbf{モザイク合成}です。\LS は、パネルの並び方（何行何列か）を指定しなくても、重なった部分の模様から
互いの位置を求めてつなぎます。

\begin{point}[2通りの使い方]
  \begin{itemize}[itemsep=1pt, topsep=2pt]
    \item \textbf{パネルごとの動画から}：キューに入れた動画を1本ずつスタックし、終わったら自動でつなぎます（第\ref{sec:mosaic-queue}節）。
    \item \textbf{スタック済みの画像から}：すでにスタックや仕上げを済ませた画像（TIFF など）を選んでつなぎます（第\ref{sec:mosaic-files}節）。
  \end{itemize}
  どちらも、つないだ結果は「仕上げ・出力」タブに送られ、ウェーブレットや色の調整をしてから書き出せます。
\end{point}

\begin{tip}[撮るときのこつ]
  \begin{itemize}[itemsep=1pt, topsep=2pt]
    \item 隣り合うパネルは\textbf{2〜3割}重ねます。重なりが少ないと、位置を求める手がかりが足りなくなります。
    \item 露出・ゲインは全部のパネルで同じにします。薄雲などで明るさが少し変わるのは、\LS がそろえます。
    \item 月は少しずつ欠け方や向きが変わるので、全部のパネルを続けて短い時間で撮ります。
  \end{itemize}
\end{tip}

% ------------------------------------------------------------------------------
\section{パネルごとの動画をまとめて処理する}\label{sec:mosaic-queue}

\begin{step}{1}{パネルの動画を全部キューに入れる}
  パネルごとの動画（または静止画連番のフォルダ）を、入力キューへまとめてドラッグ＆ドロップします。並べる順番は問いません。
\end{step}

\begin{step}{2}{\uil{モザイク合成} のメニューで \uil{キューの素材をすべてスタックしてつなぐ} を選ぶ}
  入力キューの下の \uil{モザイク合成} を押して選びます。\textbf{選んだだけでは処理は始まりません}。
  画面下のボタンが \ui{現在の設定でモザイク合成を含めた一括処理} に替わります。
\end{step}

\screen{mosaic_ready}{モザイク合成を選んだところ（まだ実行していない）}{%
  \mk{queue}{1}\mk[e]{mosaicmenu}{2}\mk[n]{runall}{3}\mk[nw]{stages}{4}\mk[w]{tabs}{5}\mk[w]{status}{6}}

\begin{marklist}
  \item キューに入れたパネルの動画（この例は月の上半分と下半分の2本）。
  \item \uil{モザイク合成} のメニュー。選んである間は、メニューの中の項目に印（✓）が付きます。\textbf{もう一度選ぶと元に戻ります}。
  \item \ui{現在の設定でモザイク合成を含めた一括処理}：これを押すと処理が始まります。
  \item この状態では、\ui{品質評価}\ui{アライメント}\ui{スタック} は押せません（どのパネルも同じ設定で通して処理するためです）。
  \item 設定は各タブでふだんどおり変えられます。\textbf{ここで決めた設定が、すべてのパネルに使われます}。
  \item 状態表示に、次にすることが出ます。
\end{marklist}

\begin{step}{3}{設定を確かめてから \ui{現在の設定でモザイク合成を含めた一括処理} を押す}
  キューの動画を上から順に、品質評価 → アライメント → スタックと進め、全部終わるとつないで、結果を入力キューに足します。
  処理中は状態表示に \uil{［パネル 1/2］} のように、いま何本目かが出ます。\ui{中断}（\key{⌘}\key{.}）で止められます。
\end{step}

\screen{mosaic_done}{モザイク合成が終わったところ（背景を黒にして、色を自動で合わせたあと）}{%
  \mk{queue}{1}\mkd[inw]{preview}{2}\mk[w]{background}{3}\mk[n]{export}{4}}

\begin{marklist}
  \item つないだ結果が、\uil{モザイク}という印の付いた行として入力キューに足されます。副題には大きさと、つなぎ目のずれ（画素）が出ます。
        元の動画の行は \uil{完了} になっています。
  \item つないだ画像。「仕上げ・出力」タブで開いているので、スタックの結果と同じように仕上げができます。
  \item 仕上げの設定。この例では \ui{背景を黒にする}（第\ref{sec:background}節）と、色の \ui{自動（灰色仮説）} を掛けています。
  \item \ui{書き出し…} で保存します。ファイル名は「最初のパネルの名前\texttt{\_mosaic}」になります。
\end{marklist}

つなぎ終わった直後の状態表示には、\uil{モザイク合成が完了しました — 2 枚中 2 枚を配置（3921×3573、つなぎ目のずれ 0.48 px）}のように、
つないだ枚数・大きさ・つなぎ目のずれが出ます。

\begin{tip}[設定を変えてやり直す]
  \uil{キューの素材をすべてスタックしてつなぐ} は、処理が終わっても選んだままになっています。採用率などを変えてもう一度ボタンを押せば、
  すぐにやり直せます（設定が変わっていない工程は飛ばすので、スタックだけを変えたなら短い時間で済みます）。
  やり直すたびに、結果は新しい行として足されます。パネルを1本ずつ処理する使い方に戻るには、メニューで同じ項目をもう一度選びます。
\end{tip}

\begin{caution}[フレームの範囲と処理範囲はパネルごと]
  \uil{使うフレームの範囲}と\uil{処理範囲}は、動画ごとに決めたものを使います。それ以外の設定は、ボタンを押したときのものが全部のパネルに使われます。
  採用率を決めるのに迷ったら、先にパネルを1本だけ選び、\ui{採用率の目安を測る}（第\ref{sec:adoption}節）で確かめておくとよいでしょう。
\end{caution}

% ------------------------------------------------------------------------------
\section{スタック済みの画像をつなぐ}\label{sec:mosaic-files}

パネルごとにスタックと仕上げを済ませて書き出した画像があるときは、それをつなげます。

\begin{enumerate}[itemsep=2pt]
  \item \uil{モザイク合成} のメニューで \uil{画像ファイル（パネル）をつなぐ…} を選びます（メニューバーの「ファイル → モザイク合成（画像ファイルから）…」でも同じ）。
  \item パネルの画像（TIFF・PNG・FITS・JPEG）を\textbf{2枚以上}選んで開きます。
  \item すぐにつなぎ、結果を「仕上げ・出力」タブで開きます。
\end{enumerate}

Finder から画像を\textbf{2枚以上まとめて}入力キューへドロップしたときは、開き方を尋ねる画面が出ます。

\begin{center}\small
\begin{tabularx}{0.95\linewidth}{@{}>{\sffamily\bfseries}l X@{}}
\toprule
連番としてスタック & 撮影したコマ（連写した写真）として、1本の連番にまとめます（第\ref{sec:raw}節）。\\
それぞれ仕上げで開く & スタック済みの画像として、1枚ずつ「仕上げ・出力」タブで開けるようにします（第\ref{sec:open-image}節）。\\
モザイク合成 & パネルとしてつなぎます。\\
\bottomrule
\end{tabularx}
\end{center}

41枚以上のとき、カメラの RAW が混ざっているときは、尋ねずに連番として追加します。

\begin{tip}[パネルごとに仕上げてからつなぐか、つないでから仕上げるか]
  どちらでもかまいません。パネルごとにウェーブレットの強さが違うと継ぎ目で見え方が変わるので、迷ったら\textbf{先につないで、あとでまとめて仕上げる}
  （前の節の方法）ほうが、全体がそろった仕上がりになります。
\end{tip}

% ------------------------------------------------------------------------------
\section{つなぎ方と、うまくいかないとき}

\LS は次のことを自動で行います。設定する項目はありません。

\begin{itemize}[itemsep=2pt]
  \item \textbf{位置}：重なった部分の模様から、パネルどうしの位置を1画素より細かく求めます。
  \item \textbf{向き}：経緯台で撮ると、時間とともに視野が少しずつ回ります。小さな回転（数度まで）と倍率のわずかな差は補正します。
  \item \textbf{明るさ}：パネルごとの明るさの差をそろえます。空が写っているときは、空の明るさ（カメラの下駄）より上にだけ倍率を掛けるので、
        空の明るさが継ぎ目で変わることはありません。
  \item \textbf{継ぎ目}：なだらかな明るさは広い範囲で混ぜ、細かい模様は継ぎ目の近くだけで切り替えます（細部が二重になりません）。
\end{itemize}

\begin{tabularx}{\linewidth}{@{}>{\raggedright\arraybackslash}p{0.36\linewidth} X@{}}
\toprule
こんなとき & どうするか\\
\midrule
\uil{パネルどうしの重なりを見つけられませんでした} & 隣り合うパネルが十分に重なっているか（2〜3割）確かめます。模様の乏しい所（海の中央など）だけで重なっていると、位置が決まりません。\\
\uil{配置できなかったパネルがあります} & ほかのどのパネルとも重なりが見つからなかったパネルは、入れずにつなぎます。そのパネルの名前が出るので、重なりを確かめてください。\\
つないだ画像の隅が黒い & どのパネルも届かなかった所は黒のままになります。「向き・切り抜き」で切り抜いてください（第\ref{sec:crop}節）。\\
空が黒くならない・色が付く & \ui{背景を黒にする} を押します（第\ref{sec:background}節）。\\
\bottomrule
\end{tabularx}
